โครงงานนี้มีวัตถุประสงค์เพื่อศึกษา ออกแบบ และพัฒนาระบบจำแนกและวิเคราะห์ข่าวสารอัตโนมัติด้วยเทคนิคการประมวลผลภาษาธรรมชาติแบบเฉพาะประเภท (Category-Specific News Classification and NLP Analytics Pipeline) ที่ครอบคลุมการทำงานตั้งแต่การรวบรวมข่าวสารจากหลายสำนักข่าว การสกัดเนื้อหา การสรุปความ การจำแนกหมวดหมู่ การวิเคราะห์แนวโน้มความนิยม และการนำเสนอข้อมูลแบบเรียลไทม์ โดยสถาปัตยกรรมระบบได้รับการออกแบบตามหลักการ SOLID และ GRASP แบ่งการทำงานเป็นเลเยอร์อย่างเป็นระเบียบ ส่วนบริการเบื้องหลัง (Backend) พัฒนาด้วยภาษา Python และ FastAPI รองรับการดึงข้อมูลทั้งจาก RSS Feed และการสกัดหน้าเว็บผ่าน Playwright Headless Browser พร้อมระบบแคชแบบหลายระดับ (Multi-tier Caching: หน่วยความจำและ HTTP 304) เพื่อลดภาระการเชื่อมต่อ การสรุปเนื้อหาและวิเคราะห์อารมณ์ของข่าว (Sentiment) ทำงานร่วมกับโมเดลภาษาขนาดใหญ่ (LLM) ผ่าน KKU Gen.AI API ขณะที่การจำแนกหมวดหมู่ข่าวภาษาไทย 7 หมวดหมู่หลักใช้แนวทางไฮบริดผสานกฎคำสำคัญเชิงบริบท (Rule-based Cues) ร่วมกับโมเดลการเรียนรู้ของเครื่อง (SVM/TF-IDF และ WangchanBERTa) ส่วนติดต่อผู้ใช้ (Frontend) พัฒนาด้วยรูปแบบ Single Page Application (SPA) ด้วย Vanilla JavaScript แบบ ES Modules และ Tailwind CSS สื่อสารกับเซิร์ฟเวอร์แบบเรียลไทม์ผ่าน WebSocket (Socket.IO) พร้อมฟังก์ชันหน้าต่างอ่านข่าวในระบบ (In-app Reader Modal) และระบบแนะนำข่าวสารตามความสนใจเฉพาะบุคคลโดยคำนึงถึงความเป็นส่วนตัวตามมาตรฐาน PDPA จากผลการประเมินการทำงานพบว่าระบบสามารถทำงานได้อย่างมีเสถียรภาพ ตอบสนองคำขอรวดเร็ว และช่วยลดเวลาในการติดตามข่าวสารจากหลากหลายแหล่งได้อย่างมีประสิทธิภาพ